% SPDX-FileCopyrightText: 2016-2026 Lari Natri <lari.natri@iki.fi>
% SPDX-License-Identifier: GPL-3.0-or-later
%
% Song IDs, position records, anchors, prelude metadata/metrics and audio keys.
% This file is part of the 'ulsbs' package.
%
% Load: ulsbs.sty, after melody and before attached notes in the content phase.
% Requires: PDF layer setup, registered variant presentation, final typography
% and style, hyperref and keyval. The song owner consumes this state later.
% State: owns the position output stream, engine-position adapters, anchor
% deduplication, song metadata, prelude geometry and audio key values. Inherits
% @ as a letter with no enclosing group; retains local probes and global IDs.
% Begin/end-document hooks register now. Font initialization, title measurement
% and em offsets execute eagerly here in their original font context. Public
% providecommand prelude hooks can be defined before loading or renewed later,
% but later changes do not rerun measurements. Allocated lengths use setlength;
% per-song fields are populated by lifecycle keys. Audio capture preserves its
% local active-accidental catcodes and defers key assignment until use.

\ProvidesFile{ulsbs-song-metadata.code.tex}[ULSBS song metadata and prelude metrics]

% === Song IDs, physical positions and PDF anchors ===

% Machine-readable output file: <jobname>.ulsbs.psv
\newwrite\ulsbs@idtopagefile

% Cross-engine support for recording physical positions:
% - pdfTeX:   \pdfsavepos / \pdflastypos (sp)
% - LuaTeX:   \savepos    / \lastsavedpos (sp)
\newif\ifulsbs@havepos
\ulsbs@haveposfalse

\begingroup
  % Prefer pdfTeX-style primitives if available
  \ifx\pdfsavepos\undefined
  \else
    \ifx\pdflastypos\undefined
    \else
      \global\ulsbs@havepostrue
      \gdef\ulsbs@savepos{\pdfsavepos}%%
      \gdef\ulsbs@lastypos{\pdflastypos}%%
    \fi
  \fi
  % Fallback to LuaTeX primitives if pdfTeX ones are missing
  \ifulsbs@havepos\else
    \ifx\savepos\undefined
    \else
      \ifx\lastypos\undefined
      \else
        \global\ulsbs@havepostrue
        \gdef\ulsbs@savepos{\savepos}%%
        \gdef\ulsbs@lastypos{\lastypos}%%
      \fi
    \fi
  \fi
\endgroup

% Open the output file when the document begins, and write a comment line(s)
% describing content.
\AtBeginDocument{%
  \immediate\openout\ulsbs@idtopagefile=\jobname.ulsbs.psv
  \write\ulsbs@idtopagefile{%
    \string# ID-NR-PAGE-Y|<id>|<songnumber>|<page>|<y-pos(sp)>
  }%
}

% Close the output file when the document ends.
\AtEndDocument{%
  \immediate\closeout\ulsbs@idtopagefile
}

% Write one line to the psv file: ID-NR-PAGE-Y|<id>|<songnumber>|<page>|<y-pos>
% Note that \ulsbs@anchorbyid should be called before this command, as it
% will create a deduplicated form of the song's id this macro uses if
% available.
\newcommand\ulsbs@writeidtopage{%
  \ifx\ulSongId\@empty
    % no id for this song -> write nothing
  \else
    \begingroup
      % Use the resolved id (with possible -2, -3, ...) if available;
      % fall back to \ulSongId if for some reason it wasn't set.
      \edef\ulsbs@fileid{%
        \ifx\ulsbs@idwithsuffix\relax
          \ulSongId
        \else
          \ulsbs@idwithsuffix
        \fi
      }%
      % Freeze id and songnum *now*; delay \thepage (and y-pos) to shipout
      \edef\ulsbs@tmp{%
        \write\ulsbs@idtopagefile{%
          ID-NR-PAGE-Y|\ulsbs@fileid|\theulSongNumber|\noexpand\thepage%
          \ifulsbs@havepos
            |\noexpand\the\ulsbs@lastypos%
          \fi
        }%
      }%
      % Asynchronous write: expands \thepage and y-pos at shipout
      \ulsbs@tmp
    \endgroup
  \fi
}

% Invisible anchor for each song, based on \ulSongId
% If the same \ulSongId occurs multiple times, the first gets
%   songid:<id>
% and later ones get
%   songid:<id>-2, songid:<id>-3, ...
% This deduplication is done for the anchors only, not for the PSV file
% written by \ulsbs@writeidtopage
\newcommand\ulsbs@anchorbyid{%
  \ifx\ulSongId\@empty
    % no id for this song -> no anchor, no file id
  \else
    % Always compute a unique id (even without hyperref), and store it in
    % \ulsbs@idwithsuffix (e.g., 'forca-da-cura-a4fe5a', 'forca-da-cura-a4fe5a-2', ...)
    % and \ulsbs@anchorname (e.g., 'songid:<id[-n]>').
    \begingroup
      \edef\ulsbs@id{\ulSongId}%
      % Check whether we've seen this id before
      \expandafter\ifx\csname ulsbs@seen@\ulsbs@id\endcsname\relax
        % First occurrence: remember it as 1, no suffix
        \expandafter\xdef\csname ulsbs@seen@\ulsbs@id\endcsname{1}%
        \xdef\ulsbs@idwithsuffix{\ulsbs@id}%
      \else
        % Later occurrence: increment counter and add -<n> suffix
        \expandafter\xdef\csname ulsbs@seen@\ulsbs@id\endcsname{%
          \number\numexpr\csname ulsbs@seen@\ulsbs@id\endcsname+1\relax}%
        \xdef\ulsbs@idwithsuffix{\ulsbs@id-\csname ulsbs@seen@\ulsbs@id\endcsname}%
      \fi
      \xdef\ulsbs@anchorname{songid:\ulsbs@idwithsuffix}%
    \endgroup
    % Only if hyperref's hypertarget is available, create the actual anchor
    \ifx\hypertarget\undefined\else
      \ifx\hypertarget\relax\else
        % Record current position (for y-offset) if available
        \ifulsbs@havepos
          \ulsbs@savepos
        \fi
        \hypertarget{\ulsbs@anchorname}{}%
      \fi
    \fi
  \fi
}

% The song lifecycle calls the anchor before writing its page record.

% To set this for a song, use the optional id={} parameter for \beginsong.
% NOTE: Id's are to consist only of lowercase letters, numbers and a dash.
% Recommended format: song's main title sluggified, spaces converted to dashes,
%                     followed by 6-char random string.
% Example:
%   \beginsong{Força da Cura}[id={forca-da-cura-a4fe5a}]
%
\newcommand{\ulSongId}{}


% === Prelude metadata and measurements ===

% Font/text providecommand defaults remain beside their prelude consumers:
% define before loading or renew afterward. They expand during prelude building.
% The measurements below execute now; changing font hooks later does not repeat
% them. Preserve their order and font context even when reorganizing defaults.

% Shared row/sample font and per-field style defaults follow registered
% presentation, not visibility or mutable legacy identity. Authored hooks win.
\providecommand{\ulPreludeBadgeFont}{\ulsbs@variantbadgefont\normalfont}

% Tag-list style, applied inside the tag block; lyric presentation defaults gray.
\providecommand{\ulPreludeTagStyle}{\ulsbs@variantbadgestyle}

% Text style for the key line in the prelude, meant for changing the font
% and size -- colors, italics and bolds are set implicitly. Can be set before
% including this file.
\providecommand{\ulPreludeKeyLineStyle}{\ulsbs@variantbadgestyle}

\providecommand{\ulPreludePhaseStyle}{\ulsbs@variantbadgestyle}

% Support for "offered to" information, used in astral books. This is displayed
% after the by field in the song prelude, prepended with "offered to:" text and
% parenthesized.
% To disable displaying this extra info in the song preludes for the whole
% document, set \ulShowOfferedTo{false} in the main document preamble (after including
% this package).
% To set this for a song, use the optional off={} parameter for \beginsong.
\newcommand{\ulOfferedTo}{}


% Support for extra song information. The extra information can be used for
% example to contain information about song origins or language.
% To disable displaying this extra info in the song preludes for the whole
% document, set \ulShowExtraInfo{false} in the main document preamble (after including
% this package).
% To set this for a song, use the optional ex={} parameter for \beginsong.
\newcommand{\ulExtraInfo}{}


% Support for key information. This key is supposed to hold the musical key
% (for example "C" or "Am" without quotes) of the song, after transposing.
% So remember to change this if you transpose.
% To enable displaying these keys in the song preludes for the whole
% document, set \ulShowKey{true} in the main document preamble (after including
% this package).
% To set this for a song, use the optional key={} parameter for \beginsong.
\newcommand{\ulSongKey}{}


% Texts for key info in prelude. Can be set before including this file.
\providecommand{\ulPreludeTextKey}{\ulIfLanguage{portuguese}{Tom:}{Key:}}



\providecommand{\ulPreludeTextGoodKeysAlone}{\ulIfLanguage{portuguese}{Boas tons:}{Good keys:}}

% Support for a list of good keys (range) for a song. This keyval is supposed
% to hold the good (for singing) keys this song to be played in.
%
% To set this for a song, use the optional gk={} parameter for \beginsong.
% The preferred format for this text is {X, Y--Z}, where X is the preferred
% "best" key, and the Y--Z is the range of keys where the melody stays within
% a range limited by G note in the low end, and the B note of the next octave
% in the high end. Example: gk={Bm, Gm--Cm}. If such a range is impossible,
% put the keys in the range in parenthesis.
%
% The good keys list is displayed in the song preamble. To disable displaying
% this for the whole document, set \ulShowGoodKeys{false} in the main document
% preamble (after including this package).
\newcommand{\ulGoodKeys}{}



% Support for phase hints for songs.
% To disable displaying these phases in the song preludes for the whole
% document, set \ulShowPhase{false} in the main document preamble (after including
% this package).
% To set this for a song, use the optional ph={} parameter for \beginsong.
\newcommand{\ulPhaseList}{}%


% Tag metadata

% Support for tags for songs.
% To disable displaying the tags in the song preludes for the whole
% document, set \ulShowTags{false} in the main document preamble (after including
% this package).
% To set these for a song, use the optional tags={} parameter for \beginsong.
% An example: \beginsong{Songname}[tags={fire, water}]. Valid tags are defined
% in 'include/tags.txt', or the file selected with \ulTagFile in the preamble.

\newcommand{\ulTagList}{}% defined below for each song


% Prelude font initialization and geometry used by the song renderer.

% Initialize the number printer's fonts without retaining an unused height.
\begingroup\settoheight\@tempdima{\ulPrintSongNumber{1}}\endgroup
% Title height with default letters (no accented capitals), used in \ulMakePrelude.
\newlength\ulSongTitleDefaultHeight
\settoheight\ulSongTitleDefaultHeight{\ulSongTitleFont\relax Tih}

% Load-time offsets for the prelude's right-side metadata box, relative to the
% main-title baseline. These em dimensions capture the current font here;
% measuring the title above does not select its font for these assignments.
\newdimen{\ulsbs@preludeRTfirstraise}\setlength{\ulsbs@preludeRTfirstraise}{1.03em}%
\newdimen{\ulsbs@preludeRTsecondraise}\setlength{\ulsbs@preludeRTsecondraise}{0pt}%
\newdimen{\ulsbs@preludeRTthirdraise}\setlength{\ulsbs@preludeRTthirdraise}{-1.03em}%
\newdimen{\ulsbs@preludeRTfourthraise}\setlength{\ulsbs@preludeRTfourthraise}{-2.06em}%

% The UI point remains available without badges. These discarded boxes measure
% enabled, nonempty badges' possible top edge; the marker never reserves space.
% Hidden/absent badges must not execute their display hooks even for this sample.
\newcommand{\ulsbs@pdfheaderanchor}{%
  \ifulsbs@pdflayers\ifx\ulsbs@pdfpreludesongid\@empty\else
    \begingroup
    \dimen@=0pt
    \ifulShowKey\ifx\ulSongKey\@empty\else
      \sbox\ulsbs@pdfanchorbox{{\ulPreludeBadgeFont\ulPreludeKeyLineStyle
        \ulsbs@mainkeybox{\symudarrow\textbf{\ulsbs@vphantommaxtext}}}}%
      \dimen@=\ht\ulsbs@pdfanchorbox
    \fi\fi
    \ifulShowTags\ifx\ulTagList\@empty\else
      \sbox\ulsbs@pdfanchorbox{{\ulPreludeBadgeFont\itshape\ulPreludeTagStyle
        \ulsbs@tagsbox{\ulsbs@vphantommaxtext}}}%
      \ifdim\ht\ulsbs@pdfanchorbox>\dimen@\dimen@=\ht\ulsbs@pdfanchorbox\fi
    \fi\fi
    \ifulShowPhase\ifx\ulPhaseList\@empty\else
      \sbox\ulsbs@pdfanchorbox{{\ulPreludeBadgeFont\rmfamily\ulPreludePhaseStyle
        \ulsbs@tagsbox{\ulsbs@vphantommaxtext}}}%
      \ifdim\ht\ulsbs@pdfanchorbox>\dimen@\dimen@=\ht\ulsbs@pdfanchorbox\fi
    \fi\fi
    \makebox[0pt][l]{\raisebox{\dimexpr\ulsbs@preludeRTfirstraise+\dimen@\relax}[0pt][0pt]{%
      \kern2bp%
      \ifulsbs@havepos\ulsbs@savepos\fi
      \edef\ulsbs@pdfanchorcode{\noexpand\latelua{
        ulsbs_multikey.anchor("\ulsbs@pdfpreludesongid")}}%
      \ulsbs@pdfanchorcode
    }}%
    \endgroup
  \fi\fi
}

% === Audio metadata ===

% Audio system is used to set URL's for audio for songs and chapters. These
% will be utilized in the future.

% keyval: Setup keys to be used as \audio macro's argument:
\define@key{audio}{title}{\def\ulsbs@audio@title{#1}}%
\define@key{audio}{key}{\def\ulsbs@audio@key{#1}}%
\define@key{audio}{pitch}{\def\ulsbs@audio@pitch{#1}}%
\setkeys{audio}{title={},key={},pitch={440}}%

% Use this command to set a URL link to audio for the current song or chapter.
% Must be called within a song or a chapter. Multiple \audio definitions can be
% set for each song/chapter; the first is the primary one. Mandatory argument
% is the URL, and the preceding optional argument contains keyval pairs,
% currently supported keys are 'title', 'key' and 'pitch'. Pitch is the
% frequency of A (default: 440).
% In the URL, TeX special characters must be escaped as usual. The key value
% accepts # and & as sharp and flat accidentals.

%
% Example: \audio[title={Good version},key={C#m}]{https://link-to-audio.com}
% Delay argument collection so raw accidentals are tokenized like chord input.
\newcommand{\audio}{%
  \begingroup%
  \catcode`\#\active\catcode`\&\active%
  \ulsbs@audio%
}
\newcommand{\ulsbs@audio}[2][]{%
  \endgroup%
  \setkeys{audio}{#1}% Set new keys
  % Access title with: \ulsbs@audio@title,
  %          key with: \ulsbs@audio@key,
  %        pitch with: \ulsbs@audio@pitch,
  %          URL with: #2
}
